% Pandoc template for the fed_fraud documents.
% Based on proposal_gen/templates/detailed_proposal_template/detailed_1_classic.tex
% (Classic report): centered title page, table of contents, ruled section
% headings and the Geodesic house footer. Fonts are Latin Modern (lmodern).
% Two layouts: the full report (title page and contents) and, with
% -V brief, a compact one-page-brief header in the same style.
\documentclass[letterpaper,11pt]{article}
\usepackage[T1]{fontenc}
\usepackage[utf8]{inputenc}
\usepackage{lmodern}
\usepackage[letterpaper,top=0.9in,bottom=1.2in,left=1in,right=1in,footskip=0.6in]{geometry}
\usepackage[table]{xcolor}
\usepackage{geodesicbrand}
\usepackage{amsmath,amssymb}
\usepackage{enumitem}
\usepackage{array}
\usepackage{calc}
\usepackage{longtable}
\usepackage{makecell}
\usepackage[most]{tcolorbox}
\usepackage{titlesec}
\usepackage{needspace}
\usepackage{etoolbox}
\usepackage{hyperref}
\hypersetup{colorlinks=true,urlcolor=coolblack,linkcolor=coolblack,filecolor=geoblue,
  pdftitle={Recovery and ROI model},pdfauthor={\geofirm}}
\urlstyle{same}

\setlist{itemsep=2pt,leftmargin=1.4em}
\setlength{\parindent}{0pt}
\setlength{\parskip}{6pt}
\providecommand{\tightlist}{\setlength{\itemsep}{0pt}\setlength{\parskip}{0pt}}
\setcounter{secnumdepth}{0}

% Section style: house rule under each heading
\titleformat{\section}{\Large\bfseries\color{coolblack}}{\thesection}{0.6em}{}[{\color{coolblack}\titlerule[0.5pt]}]
\titleformat{\subsection}{\large\bfseries\color{geomid}}{\thesubsection}{0.6em}{}
\titleformat{\subsubsection}{\normalsize\bfseries\color{geomid}}{\thesubsubsection}{0.6em}{}
\titlespacing*{\section}{0pt}{10pt}{6pt}

% Tables are written by latex/filters.lua (navy header, banded rows)
\setlength{\LTpre}{6pt}\setlength{\LTpost}{6pt}

\begin{document}

% BRIEF HEADER
\thispagestyle{geodesic}
\begin{minipage}[c]{0.13\textwidth}\includegraphics[width=\linewidth]{icon.png}\end{minipage}\hfill
\begin{minipage}[c]{0.84\textwidth}
  {\textsc{\large \geofirm}}\par\vspace{2pt}
  {\small\color{geogray}Federal fraud expansion\enspace\textperiodcentered\enspace When: \compiletime}
\end{minipage}\par\vspace{6pt}
{\color{coolblack}\rule{\linewidth}{0.3mm}}\par\vspace{2pt}
{\LARGE\bfseries\color{coolblack} Recovery and ROI model}\par
\vspace{2pt}{\large\color{geomid} The business case by program}\par
\vspace{-2pt}{\color{coolblack}\rule{\linewidth}{0.3mm}}\par\vspace{4pt}
\pagestyle{geodesic}

\textbf{Status:} Draft v0.1 (2026-11-04), for Steve's review. Extension
3 of 10 (see the \href{../extensions.pdf}{extensions overview}). Model:
\texttt{analytics/roi/model.py}; inputs:
\texttt{analytics/roi/programs.csv}; results:
\texttt{docs/proposals/data/roi\_results.csv}. Interactive version on
the site at \texttt{/tools/roi/}. \textbf{These are estimates for a
business case, not savings. Savings stay Pending until an Inspector
General or DOJ confirms them.}

\hypertarget{bottom-line}{%
\section*{Bottom line}\label{bottom-line}}
\addcontentsline{toc}{section}{Bottom line}

Leaders fund what pays back. Starting from each program's reported
improper payments (GAO, FY2025) and a stated chain of assumptions, the
model estimates the five-year present value of what detection tooling
would prevent or recover, against what it costs to build, run and
review.

\begin{itemize}
\tightlist
\item
  \textbf{The large health and benefit programs pay back many times over
  even on cautious assumptions.} Medicare fee-for-service: median net
  benefit about \$880M over five years, roughly 10 times cost; Medicaid
  about \$450M (8 times), even with only a fifth of Medicaid reachable
  through state data.
\item
  \textbf{The break-even bar is low.} For Medicare, the tooling pays for
  itself if it catches 0.2\% of the improper dollars it can see that
  current controls miss. For SNAP it is 0.8\%. A pilot only has to beat
  that number.
\item
  \textbf{Small programs do not pay back on their own.} OPM retirement,
  FEMA Public Assistance and the closed Shuttered Venue grant program
  have a 50\% or lower chance of positive net benefit as stand-alone
  builds. They should ride on shared screens (cross-agency matching,
  exclusion checks), not get their own tooling.
\end{itemize}

\hypertarget{results-by-program}{%
\section*{Results by program}\label{results-by-program}}
\addcontentsline{toc}{section}{Results by program}

\begingroup\small\setlength{\tabcolsep}{4pt}\renewcommand{\arraystretch}{1.25}
\rowcolors{2}{white}{geolight}
\begin{longtable}[l]{@{}>{\raggedright\arraybackslash}p{\dimexpr 0.282\linewidth-2\tabcolsep}l>{\raggedright\arraybackslash}p{\dimexpr 0.197\linewidth-2\tabcolsep}>{\raggedright\arraybackslash}p{\dimexpr 0.197\linewidth-2\tabcolsep}ll@{}}
\rowcolor{coolblack}\leavevmode\color{white}\bfseries Program & \leavevmode\color{white}\bfseries Improper & \leavevmode\color{white}\bfseries Benefit / cost & \leavevmode\color{white}\bfseries Net, P10 to P50 & \leavevmode\color{white}\bfseries ROI & \leavevmode\color{white}\bfseries Break-even \\
\endhead
Medicare FFS & \$28.8B & \$993M / \$100M & \$499M to \$884M & 9.7x & 0.17\% \\
Medicaid & \$37.4B & \$516M / \$62M & \$249M to \$449M & 8.0x & 0.32\% \\
EITC & \$21.1B & \$291M / \$45M & \$131M to \$244M & 6.4x & 0.57\% \\
SSI & \$7.3B & \$252M / \$42M & \$111M to \$208M & 5.9x & 0.66\% \\
SNAP & \$10.2B & \$211M / \$38M & \$89M to \$171M & 5.4x & 0.79\% \\
VA (7 programs) & \$3.9B & \$135M / \$32M & \$49M to \$102M & 4.2x & 1.2\% \\
HUD rental aid & \$4.3B & \$119M / \$31M & \$41M to \$87M & 3.8x & 1.4\% \\
UI & \$5.6B & \$116M / \$31M & \$39M to \$85M & 3.8x & 1.4\% \\
SSDI & \$2.3B & \$79M / \$28M & \$20M to \$51M & 2.9x & 2.1\% \\
SVOG (closed) & \$10.1B & \$32M / \$31M & -\$14M to \$1M & 1.0x & 6.9\% \\
OPM retirement & \$0.4B & \$20M / \$23M & -\$13M to -\$3M & 0.9x & 8.3\% \\
FEMA PA & \$0.3B & \$12M / \$22M & -\$18M to -\$9M & 0.6x & 13\% \\
\end{longtable}\endgroup

Improper payments are FY2025. Benefit, cost and ROI are five-year
medians; break-even is the catch rate at which net benefit is zero.
Figures are present values at a 2\% real discount rate, from 20,000
Monte Carlo draws. Run \texttt{python3\ analytics/roi/model.py} to
reproduce them.

\hypertarget{the-chain}{%
\section*{1. The chain}\label{the-chain}}
\addcontentsline{toc}{section}{1. The chain}

For each program and year:

\begin{itemize}
\tightlist
\item
  \textbf{caught} = improper payments x reach x detectable x new
  detection x ramp
\item
  \textbf{flagged} = caught / precision (what reviewers have to look at)
\item
  \textbf{benefit} = caught x {[}prevent + (1 - prevent) x recovery{]} x
  (1 + deterrence)
\item
  \textbf{cost} = build (year 1) + run (every year) + review cost x
  flagged
\end{itemize}

Closed programs get no prevention and no deterrence: only recovery on
the stock of improper payments, spread over three years.

\hypertarget{assumptions-and-why}{%
\section*{2. Assumptions and why}\label{assumptions-and-why}}
\addcontentsline{toc}{section}{2. Assumptions and why}

Each share is drawn from a triangular distribution (low, likely, high).
These are the numbers to argue about with an agency, and the pilot is
designed to measure the three that matter most (new detection,
precision, prevent).

\begingroup\small\setlength{\tabcolsep}{4pt}\renewcommand{\arraystretch}{1.25}
\rowcolors{2}{white}{geolight}
\begin{longtable}[l]{@{}>{\raggedright\arraybackslash}p{\dimexpr 0.159\linewidth-2\tabcolsep}>{\raggedright\arraybackslash}p{\dimexpr 0.207\linewidth-2\tabcolsep}ll>{\raggedright\arraybackslash}p{\dimexpr 0.478\linewidth-2\tabcolsep}@{}}
\rowcolor{coolblack}\leavevmode\color{white}\bfseries Assumption & \leavevmode\color{white}\bfseries Low & \leavevmode\color{white}\bfseries Likely & \leavevmode\color{white}\bfseries High & \leavevmode\color{white}\bfseries Basis \\
\endhead
Reach & per program (0.2 to 1.0) &  &  & Share of program dollars our data access covers within five years. Medicaid and EITC 0.2 (state data, tax-data limits); Medicare 0.5; OPM 0.8 \\
Detectable & 25\% & 40\% & 55\% & Share of improper dollars with a data trace (identity, eligibility, billing pattern) rather than a missing document. Agencies report most improper payments as documentation errors; the detectable share is what the screens can reach \\
New detection & 3\% & 8\% & 15\% & Share of detectable dollars our screens catch that current controls miss. The single most uncertain number; the pilot measures it \\
Precision & 30\% & 50\% & 70\% & Share of flagged dollars confirmed on review. Sets reviewer workload \\
Prevent & 20\% & 45\% & 70\% & Share caught before payment. Needs pre-payment integration (Do Not Pay, claims edits) \\
Recovery & 5\% & 15\% & 30\% & Share recovered when caught after payment. Pay-and-chase recovers little; GAO and OIG reports show low recovery rates \\
Deterrence & 0\% & 5\% & 15\% & Extra benefit from fewer attempts once controls are known \\
Build & \$2M & \$4M & \$8M & One-time, per program, on the shared platform \\
Run & \$1.5M & \$3M & \$6M & Per year, per program: data, compute, analysts \\
Review cost & 1\% & 2\% & 4\% & Review and investigation cost per flagged dollar \\
Ramp & 25\%, 60\%, 100\%, 100\%, 100\% &  &  & Capacity reached in years 1 to 5 \\
Discount rate & 2\% real &  &  & OMB Circular A-94 (2023 revision) default; check the current value \\
\end{longtable}\endgroup

\hypertarget{what-the-model-does-not-count}{%
\section*{3. What the model does not
count}\label{what-the-model-does-not-count}}
\addcontentsline{toc}{section}{3. What the model does not count}

\begin{itemize}
\tightlist
\item
  Criminal fines, False Claims Act damages and penalties (often larger
  than the recovered overpayment; treble damages under 31 U.S.C. 3729).
\item
  Fraud outside improper payment estimates. GAO estimated total fraud
  losses at \$233B to \$521B a year (GAO-24-105833), well above reported
  improper payments, so starting from improper payments is conservative.
\item
  Shared-platform savings: the second and later programs reuse the
  entity spine and screens, so build cost per program falls.
\item
  Benefits to states (the state share of Medicaid, UI administration).
\end{itemize}

\hypertarget{sensitivity}{%
\section*{4. Sensitivity}\label{sensitivity}}
\addcontentsline{toc}{section}{4. Sensitivity}

The ranking is driven by improper dollars x reach. Within a program, net
benefit moves almost one for one with new detection and prevent.
Precision matters through review cost only; at 30\% precision review
cost roughly doubles, but no large program turns negative. The
break-even column is the best single number for a leader: if the pilot
shows a catch rate above it, scale; if not, stop.

\hypertarget{how-to-use-it-in-a-proposal}{%
\section*{5. How to use it in a
proposal}\label{how-to-use-it-in-a-proposal}}
\addcontentsline{toc}{section}{5. How to use it in a proposal}

\begin{enumerate}
\def\labelenumi{\arabic{enumi}.}
\tightlist
\item
  Quote the program's row, the break-even rate and the P10 (pessimistic)
  net, not the median.
\item
  Replace our assumptions with the agency's own where it has them
  (recovery rates, review cost).
\item
  State that savings stay Pending until confirmed, and that the pilot
  exists to measure the uncertain assumptions.
\end{enumerate}

\end{document}
